\bookmark[dest=\HyperLocalCurrentHref,level=1]{New NLOS Scenes}
\section{New NLOS Scenes} \label{sec5}
In Sec.~\ref{sec2}-\ref{sec4}, most existing NLOS imaging studies have been introduced. These studies are accomplished in active and passive NLOS scenarios, as shown in Fig.~\ref{fig:introFig}. However, NLOS imaging is not limited to these scenes, and more new types of scenes are waiting to be developed. In this section, we take the recent two bounce NLOS imaging~\cite{vedaldi_imaging_2020} and keyhole NLOS imaging~\cite{metzler_keyhole_2021} as examples to introduce potential new NLOS scenes and their prospects. We further summarize the 2022--2026 expansion of NLOS imaging into sparse and irregular active acquisition, active corner cameras, ptychographic depth-resolved NLOS, transient light-transport-matrix iteration, phasor-field enhancement under partial sampling, scattering-media robustness, few-prior passive transport estimation, low-cost LiDAR spatial reasoning, robotic navigation and exploration, radar/RF/mmWave sensing, passive electromagnetic-skin-aided radar imaging, acoustic and ultrasound sensing, terahertz waves, and semantic hidden-scene understanding.

\bookmark[dest=\HyperLocalCurrentHref,level=2]{Two-bounce NLOS imaging}
\subsection{Two bounce NLOS imaging}
Henley \etal~introduced the novel "two-bounce NLOS imaging", as shown in Fig.~\ref{fig:newNLOS}-(a)~\cite{vedaldi_imaging_2020}. In this new scene, the hidden object is blocked by an occluder, but there are two diffuse reflection surfaces $W_1$ and $W_2$ on both sides of the hidden object. The observer can reconstruct the hidden scene by scanning the laser on one of the diffuse reflection surfaces (e.g., $W_1$) and using a traditional camera to capture the two-bounce light on the other reflection surface. In other NLOS imaging, hidden objects play a role in reflecting light. However, in two bounce NLOS imaging, hidden objects play a role in blocking light. By capturing the projection of light on $W_2$ after being blocked by a hidden object, whether there is a hidden object point in each voxel can be judged.

As discussed in~\cite{vedaldi_imaging_2020}, two bounce NLOS imaging can be applied to seeing behind trucks for autonomous vehicles and imaging between windows for search and rescue. These applications are extensions to traditional NLOS imaging (i.e., three-bounce NLOS imaging or passive NLOS imaging) applications and illustrate the broad application prospects of ``turning walls into mirrors''.

\bookmark[dest=\HyperLocalCurrentHref,level=2]{Two-Bounce NLOS: Role of Transients}
\subsection{Two-Bounce NLOS: Role of Transients}
Building on the two-bounce NLOS setup, Somasundaram~\etal~provided the first systematic analysis of the role of time-of-flight (transient) information in 2B-NLOS imaging under multiplexed illumination~\cite{somasundaramRoleTransients2023}. Their study quantifies the tradeoffs between temporal resolution, spatial resolution, and number of image captures, deriving formal mathematical constraints for 2B lidar and showing that transient information can significantly reduce the measurement requirements for shape reconstruction. This analytical framework provides design guidelines for future NLOS systems that incorporate time-gating in two-bounce configurations.

A complementary steady-state trajectory replaces binary shadow carving with a continuous learned scene model. Zhang~\etal~introduced Neural Illumination Fields (NIF)~\cite{zhangNeuralIlluminationFields2026}, which jointly represents hidden-space density and illumination-dependent image formation with multilayer perceptrons and optimizes the representation through differentiable intensity rendering. Because NIF synthesizes measured shadow images directly, it avoids error-prone binary shadow segmentation and remains effective under low shadow contrast and ambient-light interference. This work shifts two-bounce NLOS from discrete occupancy carving toward self-supervised neural inverse rendering for large hidden volumes.

\vspace{0.8mm}
\noindent \textbf{Dynamic neural shadow fields for moving two-bounce scenes.}
The static-scene assumption remains restrictive when hidden people or objects move during a cyclic multi-position scan. Zhang~\etal~addressed this regime with D-NeSF~\cite{zhangDynamicNeuralShadowFields2026}, decomposing the four-dimensional hidden volume into six spatiotemporal feature planes whose interpolated features are mapped by an MLP to time-varying voxel density. A cyclic acquisition strategy aligns repeated illumination positions with periodic shadow observations, while a cumulative-transmittance constraint enables self-supervised fitting to captured binary shadow sequences. This work extends the two-bounce neural-field trajectory from high-fidelity static geometry to temporally consistent reconstruction of moving hidden targets, and exposes motion modeling and scan-time synchronization as central design variables for dynamic shadow-based NLOS.

\bookmark[dest=\HyperLocalCurrentHref,level=2]{Keyhole Imaging}
\subsection{Keyhole Imaging}
For all the NLOS imaging systems introduced above, there is a limitation that a large illumination/imaging area is required.
However, in some NLOS scenes (e.g., hidden objects are blocked in a room with curtains), large imaging areas are not allowed, and only a few small holes can be used.
Metzler~\etal~used a beam splitter to propose an imaging system similar to confocal NLOS but does not require scanning, called keyhole imaging, as shown in Fig.~\ref{fig:newNLOS}-(b)~\cite{metzler_keyhole_2021}. In this system, the pulsed laser and the time-resolved detector only use a small hole to illuminate and detect the hidden space. When the object in the hidden space moves with time, expectation-maximization can be used to complete the task of NLOS shape reconstruction and localization.

\bookmark[dest=\HyperLocalCurrentHref,level=2]{Active Corner Cameras}
\subsection{Active Corner Cameras}
A newer line of work relaxes the need for dense multi-position illumination by turning a visible corner into a structured sensing device. \href{https://www.nature.com/articles/s41467-023-39327-2}{Seidel~\etal} demonstrated an active corner camera that models the complete optical response of a visible corner and reconstructs both moving foreground objects and stationary hidden background geometry from NLOS snapshots. This direction is important because it moves NLOS imaging closer to deployed corner-mounted sensing: the system no longer only recovers a single isolated object, but can also maintain a background map of the hidden space while monitoring dynamic foreground changes.


\bookmark[dest=\HyperLocalCurrentHref,level=2]{Neuromorphic NLOS Tracking}
\subsection{Neuromorphic NLOS Tracking}
Dynamic hidden-scene sensing does not always require reconstructing a complete albedo or geometry volume. Zhu~\etal~introduced computational neuromorphic NLOS tracking, using an event camera to respond selectively to motion-induced changes in the dynamic speckle field~\cite{zhuEfficientNeuromorphicNLOSTracking2024}. Because static relay-surface and background contributions generate few events, the event stream provides a compact, non-redundant representation for direct motion estimation. This work complements frame-based passive event-camera reconstruction and transient NLOS videography by establishing an event-driven task branch focused on efficient hidden-target tracking rather than full-scene recovery.

\bookmark[dest=\HyperLocalCurrentHref,level=2]{Ptychographic NLOS Imaging}
\subsection{Ptychographic NLOS Imaging}
Most transient NLOS methods treat the relay wall as an approximately diffuse surface and rely on time-of-flight inversion. Song~\etal~introduced ptychographic NLOS imaging~\cite{songPtychographicNLOS2024}, which instead exploits coherent diffraction and coded ptychography for depth-resolved visualization of hidden objects. In this formulation, a laser spot is scanned on the relay wall, the returning object wavefields are modulated by the complex-valued wall profile, and the resulting diffraction patterns form a ptychogram. Recovering both the hidden objects at different depths and the wall modulation profile connects NLOS imaging to computational wavefront retrieval, suggesting a complementary route for high-resolution depth-aware hidden-scene imaging when coherent measurements and wall-profile modulation are available.

\bookmark[dest=\HyperLocalCurrentHref,level=2]{Transient Light-Transport-Matrix Iteration}
\subsection{Transient Light-Transport-Matrix Iteration}
Beyond reconstructing a hidden volume from a single measured transient light transport matrix (TLTM), Sultan~\etal~proposed iterating the TLTM itself~\cite{sultanIteratingTLTM2024}. By computationally focusing the measured first-order TLTM onto surfaces in the hidden scene, the relay surface can be treated as a second-order active imaging system. This viewpoint enables hidden-scene relighting, separation of direct and indirect transport inside the hidden space, and dual-photography-style operations. It is therefore not only another reconstruction method, but also a step toward programmable transient probing of hidden scenes when dense SPAD-array measurements make fuller TLTM acquisition feasible.

\bookmark[dest=\HyperLocalCurrentHref,level=2]{Sparse and Irregular Active NLOS Acquisition}
\subsection{Sparse and Irregular Active NLOS Acquisition}
Another important practical trajectory is the relaxation of dense rectangular relay scanning. Liu~\etal~introduced a Bayesian framework for NLOS imaging with arbitrary illumination and detection patterns~\cite{liuArbitraryPatternNLOS2023}. By constructing virtual confocal signals and jointly regularizing the measured signal, virtual confocal signal, and hidden object, their CC-SOCR formulation reconstructs albedo and surface normals even when relay samples are sparse, fragmented, or geometrically irregular. A closely related few-shot extension~\cite{liuFewShotSSCR2022} further emphasizes that the hidden target, the measured transient signal, and surface-based structure priors can be regularized in mixed dimensions, enabling high-quality reconstruction from very small confocal grids such as $5\times5$ measurements. Cho~\etal~take a complementary learning-based route with LEAP~\cite{choLEAP2024}: instead of directly reconstructing the hidden volume from noisy partial transients, the network predicts clean full-aperture phasor fields from partial measurements and constrains the learned representation to the informative frequency band of phasor-field propagation. This connects sparse acquisition, denoising, and learned measurement completion, showing that practical NLOS can reduce both sampling count and scan area without changing the downstream physical propagator. Together, these works bridge classical dense ToF NLOS and the newer arbitrary-relay / mobile-capture setting: they show that practical NLOS deployment depends not only on faster solvers, but also on inverse models that remain well posed when the relay surface is incomplete, irregular, or sparsely sampled.

\vspace{0.8mm}
\noindent \textbf{Masked transient pretraining.}
Shen~\etal~introduced MARMOT, a self-supervised masked autoencoder that treats transient measurements as a reusable pretraining modality rather than training a separate reconstruction network for every NLOS task~\cite{shenMARMOT2025}. Its scanning-pattern mask removes structured subsets of relay measurements so that the visible tokens are functionally equivalent to arbitrary sparse sampling, while a Transformer encoder--decoder predicts the missing transient measurements. Pretraining on the TransVerse corpus of 500,000 synthetic 3D models allows the learned transient representation to transfer to downstream reconstruction through direct feature reuse or decoder fine-tuning. This work marks a shift from task-specific transient inversion toward foundation-style measurement priors that can support sparse acquisition, transient completion, and multiple NLOS objectives.

\bookmark[dest=\HyperLocalCurrentHref,level=2]{Arbitrary Relay Surfaces with Gaussian Transient Rendering}
\subsection{Arbitrary Relay Surfaces with Gaussian Transient Rendering}
Wang~\etal~introduced a LOS-guided pipeline that estimates visible relay geometry and represents the hidden scene with 3D Gaussian primitives optimized by a differentiable transient renderer~\cite{wangGaussianTransientRendering2026}. The method supports confocal and non-confocal measurements on spatially limited, sparse, planar, and strongly curved relay regions, then synthesizes transients on a virtual relay surface for conventional volumetric reconstruction. This work connects differentiable transient rendering with deployable arbitrary-relay capture and removes the large planar wall from the core geometric assumptions of learned NLOS inversion.

\bookmark[dest=\HyperLocalCurrentHref,level=2]{Scattering-Media NLOS Imaging}
\subsection{Scattering-Media NLOS Imaging}
Most active optical NLOS models assume that light propagates through free space and scatters only at surfaces. In their final Optics Letters study, Luesia~\etal~consider a harder regime in which the hidden scene is submerged in scattering media such as fog- or smoke-like volumetric transport~\cite{luesiaScatteringPhasor2023}. By testing phasor-field reconstructions under controlled simulated scattering parameters and real diffuser measurements, the work clarifies how virtual-wave NLOS behaves when the relay path itself contains strong volumetric scattering. This direction is important for applications such as rescue, surveillance, and biomedical imaging, where NLOS systems may need to operate in optically degraded media rather than clean laboratory air.

\bookmark[dest=\HyperLocalCurrentHref,level=2]{Passive NLOS with Limited Prior Data}
\subsection{Passive NLOS with Limited Prior Data}
Passive NLOS reconstruction is especially sensitive to changes in relay surfaces, illumination, and detection geometry because the transport matrix is usually unknown and strongly scene-dependent. Gui~\etal~proposed a High-Dimensional Projection Selection (HDPS) framework for passive NLOS imaging with arbitrary scene conditions and detection patterns using only a small amount of prior data~\cite{guiHDPSPassiveNLOS2024}. Instead of assuming that a full transport matrix is available or training a separate network for every transport condition, HDPS models the hidden target, projection, and transport relation through a high-dimensional manifold whose low-dimensional projections can be adapted to different passive NLOS settings. This work is relevant to the broader trend from calibrated laboratory passive NLOS toward deployable ordinary-camera systems that must adapt to new walls, illuminations, and detection layouts with limited calibration.

\bookmark[dest=\HyperLocalCurrentHref,level=2]{Consumer LiDAR and Low-Cost Spatial Reasoning}
\subsection{Consumer LiDAR and Low-Cost Spatial Reasoning}
Research-grade transient systems have traditionally required high-power lasers, dense calibrated scanning, and expensive timing electronics. Somasundaram~\etal~introduced a motion-induced aperture-sampling model that unifies object shape, object motion, and camera motion, enabling multi-frame NLOS reconstruction, tracking, and camera localization with smartphone-grade consumer LiDAR~\cite{somasundaramConsumerLiDAR2026}. This Nature result turns previously harmful motion and low spatial resolution into a synthetic aperture and marks a shift from laboratory NLOS hardware toward plug-and-play sensing.

Low-cost sensing also changes the role of datasets and cross-modal priors. DENALI provides space--time histograms and spatial-reasoning annotations for commodity LiDAR, establishing a CVPR benchmark for learned NLOS reasoning~\cite{behariDENALI2026}. In RF imaging, GeRaF~2.0 combines visible line-of-sight geometry outside an enclosure with radar propagation inside it, using the exterior geometry to regularize neural signed-distance reconstruction of hidden surfaces~\cite{luSeeingThroughBoxes2026}. Together these works show that deployable NLOS increasingly couples inexpensive measurements with motion, visible-scene geometry, and learned priors rather than relying only on denser transient acquisition.

\vspace{0.8mm}
\noindent \textbf{Commodity continuous-wave ToF scattering maps.}
Fang~\etal~used a low-cost continuous-wave ToF camera to acquire relay-surface intensity and depth images and proposed scalable scattering mapping (SSM) for learned decoupling of object and relay-surface scattering signatures~\cite{fangScalableScatteringTOF2025}. The study collected 210,000 depth images across polypropylene, acrylic, Lambertian, painted-door, and plaster-wall relays, and demonstrated reconstruction on laboratory objects and clothing in a corridor. Unlike picosecond transient inversion, SSM calibrates a scene-dependent degradation-to-clear mapping from commodity correlation-depth measurements, broadening deployable NLOS toward inexpensive depth cameras while making the need for coverage of new object--relay combinations explicit.

\bookmark[dest=\HyperLocalCurrentHref,level=2]{Thermal NLOS through rough relay surfaces}
\subsection{Thermal NLOS through rough relay surfaces}
Ye~\etal~extend passive long-wave-infrared NLOS beyond conveniently smooth relay materials with NLOSFormer~\cite{yeThermalNLOSFormer2026}. A physics-embedded network estimates a scene-dependent thermal convolution kernel and jointly recovers hidden appearance and relative depth, while the ThermalNLOS dataset covers rough-wall conditions and dynamic sequences. Reported operation at approximately 4~fps shows that learned thermal transport can relax idealized relay-surface assumptions while retaining physical structure, expanding NLOS toward darkness, adverse visible illumination, and practical building materials.

\bookmark[dest=\HyperLocalCurrentHref,level=2]{Radar-Based NLOS Imaging}
\subsection{Radar-Based NLOS Imaging}
Millimeter-wave (mmWave) radar offers a compelling alternative modality for NLOS sensing due to its ability to penetrate non-metallic obstacles, operate in adverse lighting conditions, and provide Doppler velocity information. Early radar NLOS works~\cite{scheinerseeingstreetcornersa} demonstrated detection and tracking of hidden vehicles using Doppler radars in automotive scenarios. More recently, Lai~\etal~proposed HoloRadar~\cite{laiHoloRadar2025}, a practical system that achieves full 3D scene reconstruction (rather than mere detection and tracking) using a single mmWave radar mounted on a mobile robot. HoloRadar employs a two-stage pipeline: the first stage generates high-resolution multi-return range images capturing both line-of-sight (LOS) and NLOS reflections, and the second stage uses a physics-guided architecture that models mirror-image ray interactions to map NLOS reflections to their true locations. HoloRadar has been validated across diverse real-world indoor environments and represents a major step toward full NLOS 3D reconstruction in uncontrolled settings. The radar approach is complementary to optical NLOS: it operates through walls and in total darkness, but at lower spatial resolution than optical methods.

\vspace{0.8mm}
\noindent \textbf{Cellular ISAC hardware for NLOS sensing.}
A complementary 5G/6G ISAC trajectory asks whether communication hardware can use multipath as an around-corner sensor rather than treating it only as a channel impairment. Tosi~\etal~first demonstrated the feasibility of NLOS target detection with a 27.4~GHz commercial mmWave ISAC proof-of-concept, including channel-state-information processing that suppresses spectral replicas caused by time-division-duplex gaps~\cite{tosiFeasibilityISACNLOS2024}. The later ICT study moved from feasibility to reliable intrusion monitoring of fully occluded moving targets, adding range--Doppler detection and probability-hypothesis-density filtering for tracking and false-alarm rejection in an industrial testbed~\cite{tosiReliableISACNLOS2026}. Together, these works connect radar NLOS with standards-compatible cellular infrastructure and show a deployment path in which communication radios become opportunistic hidden-region sensors.

A parallel RF/mmWave lineage has progressively broadened the task from around-corner localization to practical hidden-shape and full-scene reconstruction. CornerRadar learned a propagation hint map for robust indoor localization across diverse layouts~\cite{yueCornerRadar2022}, while Mosaic exploited curved and planar urban reflectors to widen automotive-radar blind-spot coverage~\cite{woodfordMosaic2022}. RFlect then used practical environmental reflectors, including poles and curved/composite surfaces, for around-corner shape imaging~\cite{doddsRFlect2024}. The geometry branch advanced with mmNorm, which estimates mmWave surface normals and integrates them into hidden 3D shape~\cite{doddsMmNorm2025}, and with Jeon~\etal, who infer T-junction layout from static radar points and ray-trace dynamic multipath to localize multiple hidden pedestrians in measured outdoor scenes~\cite{jeonTJunctionMmWaveNLOS2025}. Attention-based BiScalar-AA and its related two-stage attention network convert sparse radar points to pseudo-images for learned NLOS object detection and tracking~\cite{yuBiScalarAA2025,yuTSANNLOS2025}. Most recently, Wave-Former introduced physics-aware wireless shape completion for complete 3D geometry of fully occluded objects~\cite{doddsWaveFormer2026}, whereas RISE combines AoA/AoD multipath enhancement with hierarchical diffusion for layout reconstruction and object detection from a single static radar~\cite{zhouRISE2026}. Together, these works show a clear RF trajectory from point localization, through reflector-aware imaging and surface-normal recovery, to learned full-object and scene-level geometry.

\vspace{0.8mm}
\noindent \textbf{Rough, uncertain, and unknown relay geometry in mmWave NLOS.}
A complementary radar trajectory removes the ideal planar/specular relay assumption itself. Xu~\etal~model rough relay surfaces as stochastic microfacet scatterers and deliberately exploit the resulting multipath diversity with a stochastic-geometry-aided recovery procedure~\cite{xuRoughRelayMmWave2024}. Their later multi-angle formulation uses an automotive-squint SAR model and a double-sparse image prior to separate structure shared across heterogeneous relay orientations~\cite{xuDoubleSparseMmWave2024}; when those orientations are only approximately known, a Bayesian extension treats the relay angles as latent dictionary parameters and alternates expectation-propagation inference with parameter updates~\cite{xuBayesianMmWave2025}. Lv~\etal~push calibration further by reconstructing the relay reflector directly from mmWave measurements before around-corner human sensing, removing the dependence on LiDAR-provided reflector geometry~\cite{lvRelayReflector2025}. In parallel, Hydra explicitly models multi-bounce scattering from uncontrolled intermediate objects and localizes targets beyond the transmit field of view without prior environment knowledge~\cite{mehrotraHydra2024}. A 2026 Aerospace and Electronic Systems Magazine overview consolidates these rough-relay challenges and solutions~\cite{liuRoughRelaySurvey2026}. Taken together, this branch changes relay geometry from a fixed prerequisite into an estimated, uncertain, or even beneficial part of the inverse problem, complementing HoloRadar/RFlect-style reflector-aware reconstruction and the newer learned full-scene systems.


\vspace{0.8mm}
\noindent \textbf{Multipath-exploitation inverse scattering for RF NLOS.}
Suenobu~\etal~first applied a physical-optics linearization to a known indoor T-junction, using a numerically evaluated Green tensor to encode the wall-mediated paths between the radar aperture and hidden imaging region~\cite{suenobuPhysicalOpticsNLOS2024}. Their subsequent journal study extended this formulation to multipath-exploitation imaging of PEC objects~\cite{suenobuMultipathInverseScatteringNLOS2025}. Rather than suppressing indirect returns as ghosts, the method incorporates them into the forward operator and recovers hidden target position and planar extent. Three-dimensional full-wave simulations and measured 2~GHz anechoic-chamber experiments establish a complementary RF trajectory to mirror-symmetry backprojection and learned mmWave reconstruction: environment geometry is assumed known, while target scattering is inferred through a linearized inverse problem.

Du~\etal~extended radar NLOS from room-scale mmWave perception toward long-range X-band imaging~\cite{duXBandRadarNLOS2026}. Their system exploits the longer wavelength to make relay interactions predominantly specular and combines learned dense prediction with geometry-aware reconstruction to compensate for the low native spatial resolution of radar. Prototype experiments reconstruct hidden objects at distances up to 40~m, demonstrating that modality selection can change not only robustness but also the practical range envelope of NLOS imaging.


\vspace{0.8mm}
\noindent \textbf{Measured mmWave reconstruction from physical models to fast regularized migration.}
Wei~\etal~established a measured 77~GHz MIMO foundation for around-corner 3D radar imaging, deriving mirror-symmetry backprojection, resolution behavior, and relay-surface phase-error compensation~\cite{weiMillimeterWaveNLOS3D2022}. Cai~\etal~then introduced NSIR, which separates LOS and multipath NLOS echoes before jointly estimating hidden reflectivity, total-variation structure, and reflective-surface phase error~\cite{caiPreciseRadarNLOS2024}. RM-CSTV~\cite{liuRMCSTVRadarNLOS2024} replaced a large matrix operator with range migration while retaining complex sparsity and total-variation constraints, correcting virtual-target positions and preserving contours under aperture loss and sparse sampling. These works form a model-driven radar lineage in which environmental reflections are calibrated as an imaging aperture rather than treated only as nuisance multipath.

\vspace{0.8mm}
\noindent \textbf{Self-supervised and model-unfolded radar NLOS reconstruction.}
The radar branch subsequently adopted learned optimization without abandoning its propagation operators. ACTE-Net~\cite{caiACTENetRadarNLOS2025} combines a Gaussian-convolution artifact suppressor, an energy-enhancement operator, and holographic reconstruction in an unsupervised network validated on measured hidden targets. Cai~\etal~further exploit group invariance in an equivariant-imaging framework whose TV-constrained unfolding network and adaptive peak-convolution thresholds reconstruct 2D/3D mmWave SAR images from partial noisy echoes~\cite{caiEquivariantRadarNLOS2025}. Chen~\etal~embed a fast range-migration operator in a FISTA-derived network for measured near-field mmWave reconstruction under phase error, aperture shadowing, and multipath~\cite{chenRMOperatorMmWaveNLOS2026}; a parallel 121~GHz system replaces range migration with a compact holographic operator and demonstrates the same model-unfolding trajectory at sub-terahertz wavelengths~\cite{chenDeepUnfoldingTHzNLOS2026}. Together with HoloRadar's mobile, ray-interaction-aware scene reconstruction~\cite{laiHoloRadar2025}, this progression shifts radar NLOS from isolated target focusing toward practical full-scene geometry and adaptive learned inversion while retaining explicit propagation structure.

\vspace{0.8mm}
\noindent \textbf{Visually constrained penetrative RF neural reconstruction.}
Lu~\etal~introduced GeRaF~2.0 for joint line-of-sight and non-line-of-sight geometry reconstruction from penetrative radar measurements~\cite{luSeeingThroughBoxes2026}. Rather than optimizing the hidden signed-distance field in isolation, the method uses visually observed exterior geometry to constrain RF propagation from the visible region into an enclosed hidden volume. These line-of-sight priors stabilize the zero-level set, reduce surface ambiguity, and permit physically consistent recovery of both visible and concealed geometry. This direction complements mirror-path mmWave systems such as HoloRadar: instead of relying only on environmental specular relays, it combines penetrative RF transport with cross-modal geometric evidence to regularize a neural implicit scene representation.

Recent RF/mmWave work further expands the meaning of NLOS perception beyond monostatic reconstruction. MITO~\cite{doddsMITO2025} provides a multi-spectral mmWave dataset and simulation tool with both LOS and NLOS captures of everyday objects, lowering the barrier for learning-based mmWave NLOS perception. N$^2$LoS~\cite{shiN2LoS2025} studies tag-assisted mmWave NLOS localization, using a single backscatter tag to separate tag and environmental reflections in non-penetrable NLOS settings. Backscatter-assisted indoor NLOS positioning~\cite{ruttikBackscatterNLOS2026} extends this branch from dedicated tag modulation toward passive asynchronous backscatter devices that act as virtual anchors for corridor-constrained user tracking, showing that low-power environmental tags can turn otherwise difficult multipath into useful NLOS localization evidence. Park~\etal~use camera-derived road layout to interpret mmWave radar point clouds for NLOS pedestrian localization at urban T-junctions, connecting around-corner radar perception to autonomous-driving scene understanding~\cite{parkTjunctionPedestrian2025}. \href{https://arxiv.org/abs/2603.13981}{NLOS-aided distributed MIMO/ISAC} treats reconstructed NLOS multipath components as useful information for joint over-the-air synchronization and sparse off-grid environment imaging. \href{https://arxiv.org/abs/2603.02832}{Zhang~\etal} show that even higher-order double-bounce NLOS paths, often discarded as nuisance multipath, can improve snapshot radio SLAM and reveal landmarks that are unobservable from single-bounce paths alone. Yasmeen~\etal~experimentally used a custom one-bit RIS with a 5.5~GHz monostatic radar to steer illumination around a corridor corner and recover human walking micro-Doppler signatures~\cite{yasmeenAroundCornerRIS2024}. A follow-on dual-beam RIS radar study~\cite{yasmeenDualBeamRIS2026} examines a lower-complexity one-bit quantized RIS configuration that produces dual symmetric beams, benchmarking the resulting beam steering and radar cross-section against metal reflectors and ideal single-beam RIS baselines.

\vspace{0.8mm}
\noindent \textbf{Reconfigurable propagation for physiological NLOS sensing.}
A complementary RF branch treats the relay path as a controllable component rather than a fixed environmental reflector. Li~\etal~combine visual target localization with programmable RIS beam control so that the sensing field is redirected toward a hidden chest region, and then estimate respiration and heartbeat from the returned signal using an improved variational-mode-decomposition pipeline~\cite{liRISVitalSignNLOS2025}. Tripathy~\etal~integrate a liquid-crystal RIS with a self-injection-locked radar and experimentally demonstrate contactless vital-sign monitoring after electronically steering the sensing path into an NLOS region~\cite{tripathyLCRISVitalSign2025}. These works extend NLOS RF from geometry, localization, and micro-Doppler toward physiological sensing; they should therefore be interpreted as hidden-region semantic/biomedical sensing rather than hidden-shape reconstruction.
 Passive or static metasurfaces provide another route: Tornielli Bellini~\etal~study multi-view near-field NLOS radar imaging with non-reconfigurable electromagnetic skins~\cite{tornielliEMSkinsNLOS2024}, where a moving radar synthesizes an aperture over phase-configured passive modules that focus energy into the hidden region. Chen~\etal~further analyze RIS-aided NLOS bistatic MIMO radar with arbitrary electromagnetic-vector-sensor (EMVS) arrays~\cite{chenBeyondLambdaEMVS2026}, showing that polarization-rich EMVS measurements and PARAFAC-style factorization can resolve phase ambiguities even when sensor spacings exceed the conventional $\lambda/2$ limit. Newly identified adjacent RF works further show how engineered or naturally occurring reflectors can make hidden-space RF sensing more deployable: mmMirror~\cite{yanMmMirror2025} uses a Van-Atta-array IRS with commodity FMCW radars for device-free around-corner localization, while See and Beam~\cite{bandariSeeBeam2025} uses LiDAR-visible specular surfaces to infer NLOS mmWave reflection paths and localize NLOS users for beam steering. Together, these works suggest that future NLOS systems will not be limited to single-sensor hidden-shape reconstruction, but will also include localization, mapping, synchronization, dataset-driven representation learning, reflector-aware communication/sensing, passive backscatter infrastructure, higher-order radio-SLAM multipath exploitation, engineered static and reconfigurable metasurfaces, practical quantized RIS beam control, array-geometry design, and environment-aware radio perception.

\vspace{0.8mm}
\noindent \textbf{RIS-assisted gridless and monostatic angular sensing.}
A complementary RIS trajectory targets hidden-object direction rather than reflectivity or surface reconstruction. Yuan~\etal~formulated RIS-enabled NLOS direction finding in the covariance domain~\cite{yuanRISGridlessDoA2025}: after estimating the noise variance, atomic-norm minimization recovers a Hermitian Toeplitz representation and an ADMM solver provides gridless multi-source DoA estimates with limited receive hardware. Zhang~\etal~considered a monostatic radar--RIS--target--RIS--radar path~\cite{zhangRISMonostaticDOA2026}. A phase-codebook scans the hidden angular sector, steering-vector decoupling recovers the target outer-product matrix from the RIS-element superposition, and Root-MUSIC estimates the target angles. Both studies are numerical, so they should be read as theoretical RIS-assisted NLOS localization advances rather than measured hidden-scene imaging systems.

\vspace{0.8mm}
\noindent \textbf{Urban-intersection FMCW radar perception.}
Lin and Chen studied an application-facing 77~GHz automotive-radar setting in which specular reflections from buildings provide around-corner observations at occluded urban intersections~\cite{linDeepLearningFMCWRadar2026}. Their compact one-dimensional, AlexNet-derived regression network restores interference-corrupted chirps before standard range and angular estimation, reducing simulated severe-interference errors from $5.48$~m/$18.95^{\circ}$/$10.77^{\circ}$ to $0.56$~m/$0.46^{\circ}$/$0.73^{\circ}$ for range, angle, and azimuth deviation, respectively. Unlike measured radar geometry reconstruction systems, the study is evaluated entirely in MATLAB simulation; its value is therefore as a learned signal-restoration and autonomous-driving NLOS perception branch, with real-road multipath and material validation remaining open.

\vspace{0.8mm}
\noindent \textbf{Cyclostationary micro-Doppler detection in fully NLOS scenes.}
Liyanage~\etal~built a 2.47~GHz continuous-wave MIMO radar with one transmitter and four receivers and used spectral correlation density images of propeller micro-Doppler as inputs to a compact EfficientNet-B0 detector~\cite{liyanageSUASNLOSRadar2026}. Measurements collected in indoor and semi-urban settings used cabinets, doors, and walls to block the direct path; cyclic rotor features remained visible under multipath, and the four-channel model achieved 86.11\% overall accuracy compared with 73.77\% for a single-channel baseline. This result broadens NLOS radar from vehicles, pedestrians, and static hidden geometry toward low-altitude aerial-threat sensing. It should nevertheless be distinguished from imaging: the airframes were stationary while the propellers rotated, and the reported task was binary detection rather than localization, tracking, or hidden-shape reconstruction.


\vspace{0.8mm}
\noindent \textbf{Radar--LiDAR blind-spot fusion in open-pit mines.}
Yang~\etal~studied a deployment-oriented NLOS perception problem in which haul trucks and mining structures fully occlude rear targets from roof-mounted LiDAR~\cite{yangMiningRadarLiDARNLOS2026}. Their Blind-Spot Complementary Fusion framework calibrates 3D LiDAR and 4D mmWave radar, suppresses ground-coupled multipath artifacts, identifies physical LiDAR blind regions, and injects only radar points that pass local-density and spatial-consistency tests. Real measurements at 35~m and 70~m and a fully occluded test-field target show improved cross-modal proximity and nonzero envelope-level evidence where LiDAR alone misses the target. The proposed Volume Recovery Rate is explicitly a proxy for hidden-target existence and coarse envelope coverage, not complete shape reconstruction or semantic detection. By citing HoloRadar as the reconstruction-oriented radar milestone while pursuing training-free safety cues, this work adds a complementary branch in which NLOS sensing is evaluated by downstream blind-zone risk coverage.

\vspace{0.8mm}
\noindent \textbf{Measured multipath radar recognition and activity understanding.}
Radar NLOS can support semantic decisions even when the measurements are too sparse for hidden-shape reconstruction. Zeng~\etal~used a 15~GHz stepped-frequency system and an adaptive multi-scale residual network to combine fine high-amplitude single-path cues with global multipath structure and attention-weighted local scattering features~\cite{zengMultipathFeatureFusionNLOS2026}. Measured concrete- and wood-wall experiments distinguish pedestrians, weapon-carrying human models, quadrotor drones, and tank models with 99.6\% accuracy. Zhong~\etal~then treated separated UWB-radar propagation paths as natural multiview positives for self-supervised contrastive learning~\cite{zhongMultipathContrastiveNLOS2026}. MuPhyCoNet aligns learned features with observation- and prediction-level physical quantities, reaching 94.32\% six-action accuracy with only 10\% labels on 19,500 measured spectrograms. These studies extend around-corner RF from detection and localization toward target identity and human activity, but should remain categorized as semantic sensing rather than full 2D/3D imaging.

Tornielli Bellini~\etal~push reflector-assisted RF NLOS imaging toward communication infrastructure by embedding a monostatic hidden-scene imaging mode into the initial-access beam sweep of a next-generation base station~\cite{tornielliInitialAccessNLOS2025}. A low-cost non-reconfigurable modular reflector and imaging-specific codebook extend the standard communication beams; coherent processing across the resulting viewpoints forms an effective aperture for near-field reflectivity imaging, while a maximum-likelihood velocity estimator compensates moving targets. By deriving resolution, aperture, latency, SNR, and communication--imaging trade-offs and validating the design experimentally, this work connects earlier electromagnetic-skin NLOS prototypes to standard-aware 6G ISAC deployment.


\vspace{0.8mm}
\noindent \textbf{Tag-assisted multipath disambiguation.}
Shi~\etal~introduced N$^2$LoS, which combines a single 24~GHz radar with an active backscatter tag~\cite{shiN2LoS2025}. Hybrid frequency-hopping/direct-sequence signaling separates tag returns from environmental reflectors, while frequency--spatial MUSIC resolves additional multipath components. Experiments report roughly 10--12~cm median coordinate error at 5~m, illustrating how a lightweight tag can convert uncontrolled around-corner multipath into identifiable localization paths. The final journal record appears in IEEE Transactions on Mobile Computing, vol.~25, no.~5, pp.~6002--6016 (2026), superseding the earlier arXiv-only metadata.

\vspace{0.8mm}
\noindent \textbf{Sub-terahertz polarimetric indoor mapping.}
Alburadi~\etal~demonstrated that a 222--228~GHz polarimetric FMCW radar with a mechanically scanned fan-beam antenna can combine rapid high-resolution indoor mapping with measured around-corner imaging~\cite{alburadiPolarimetricRadarNLOS2025}. The system interprets double-bounce specular paths as hidden-scene evidence, while co- and cross-polarized channels help distinguish useful NLOS returns from direct leakage and environmental clutter. This work extends the measured mmWave lineage toward sub-terahertz carrier frequencies and emphasizes polarization as an additional physical cue for practical RF NLOS imaging.

\bookmark[dest=\HyperLocalCurrentHref,level=2]{Acoustic NLOS Imaging}
\subsection{Acoustic NLOS Imaging}
Acoustic waves provide yet another physically distinct modality for NLOS sensing. Similar to optical ToF systems, acoustic NLOS imaging exploits the time-of-flight of sound waves reflected off a relay surface. The Lindell~\etal~acoustic NLOS work~\cite{Lindell:2019:Acoustic} established that microphone arrays and audio sources can reconstruct hidden 3D scenes using the same algorithmic framework as optical NLOS. Li~\etal~translated the optical ToF and $f$--$k$ migration lineage into a coherent synthetic-aperture ultrasound system~\cite{ultrasoundNLOS2025}. A scanning emitter--receiver pair preserves acoustic phase and reconstructs multiple hidden targets and complex 3D scenes at distances up to 2\,m from the scattering surface, with measured lateral and depth resolution of approximately 1\,cm. The study also reports the NLOS modulation-transfer function and releases data and code, establishing ultrasound as a quantitatively comparable, low-power and eye-safe alternative to optical transient NLOS. Acoustic NLOS imaging is particularly suited to scenarios where electromagnetic NLOS is impractical, such as underwater or through highly opaque materials.

Boger-Lombard, Slobodkin, and Katz first showed that acoustic interferometry can retrieve effective Green functions from cross-correlations of uncontrolled broadband noise, enabling passive localization and tracking of a human hidden around a corner without controlled active probing~\cite{bogerLombardPassiveAcousticCorners2023}.

The scope of acoustic NLOS has also moved beyond relay-surface reflection. Sommer and Katz exploit knife-edge diffraction at doorways and convex corners to passively localize a hidden acoustic point source in three dimensions~\cite{sommerPassiveAcousticNLOS2026}. For a doorway, its two edges act as virtual detector arrays; for a corner, the method uses the frequency-dependent diffraction-loss signature in a manner analogous to head-related transfer functions. This relay-free setting shows that edge diffraction, rather than diffuse reflection alone, can be treated as an exploitable physical transport path for hidden-scene sensing.

\vspace{0.8mm}
\noindent \textbf{Laser--acoustic fusion for hidden-human orientation sensing.}
Do{\u{g}}an combined laser and acoustic chirp measurements at the feature level for four-class hidden-person orientation recognition~\cite{doganLaserAcousticOrientationNLOS2026}. Each sensing channel contributes 21 engineered features, yielding a 42-dimensional fused representation evaluated with classical classifiers, the proposed LAO-Net, and explainable-AI analysis under controlled NLOS experiments. This work extends acoustic NLOS from source localization and geometric reconstruction toward multimodal semantic inference; its output is an orientation class rather than a hidden image, position map, or 3D surface.

Recent acoustic work has also separated hidden-source localization from volumetric echo reconstruction. Jeon~\etal~track vehicles approaching from fully occluded road regions using an Acoustic--Spatial Pseudo-Likelihood inside a particle filter~\cite{jeonSoundVehicleNLOS2025}. Their ARIL dataset adds bird's-eye-view position ground truth and, together with OVAD, turns passive acoustic NLOS sensing into an application-oriented trajectory estimate for collision avoidance. Zhai~\etal~instead model the finite obstacle edge itself: the Biot--Tolstoy--Medwin first-order diffraction response forms a non-free-field steering vector for irregular-grid frequency-domain beamforming~\cite{zhaiIrregularGridAcousticNLOS2025}. The resulting method limits main-lobe deformation as the hidden source moves and is validated with both simulations and a 32-channel microphone-array experiment. Together, these systems extend the acoustic trajectory beyond relay-wall tomography toward passive tracking and diffraction-aware source localization in practical occluded environments.

A complementary second-order regime arises when both the direct path and first-order diffraction are unavailable. Zhai~\etal~construct a Biot--Tolstoy--Medwin second-order edge-diffraction transfer model and use it as the sensing matrix of a block-sparse inverse problem~\cite{zhaiSecondOrderAcousticNLOS2025}. Fast marginalized block sparse Bayesian learning then estimates hidden-source position and strength; simulations and a 32-channel array experiment show that higher-order diffraction, usually treated as a weak residual, can become the primary sensing path under stronger occlusion.

At the task level, Wang~\etal~combine bird's-eye-view scene geometry with time--frequency and spatiotemporal acoustic spectra for detecting vehicles that are fully occluded from the camera~\cite{wangAudioVisualNLOS2026}. Their scene-aware network uses CNN, LSTM, and Conformer modules to model local spectral structure, temporal evolution, and long-range cross-modal context, reporting 94.1\% and 97.0\% accuracy on the OVAD and AOVD datasets, respectively. Unlike acoustic echo tomography or diffraction-based source localization, this work does not reconstruct hidden geometry; it establishes a complementary semantic branch in which visible environment layout conditions the interpretation of sound propagated from an unseen traffic participant.

\vspace{0.8mm}
\noindent \textbf{Material recognition from wall-mediated acoustic echoes.}
Alaku{\c{s}} and T{\"u}rko{\u{g}}lu extend this semantic branch from hidden-agent inference to physical material recognition~\cite{alakusAcousticMaterialNLOS2026}. Their ANLOS-R acquisition places eight loudspeakers and eight microphones toward a relay wall while blocking the direct path to nine target materials, so the classifier operates on indirect chirp echoes shaped by reflection, absorption, and scattering. Classical spectral and temporal descriptors are fused with multi-scale wavelet energy and entropy into a 70-dimensional representation; recurrent models reach a reported balanced accuracy and macro-F1 of 0.99, and SHAP analysis relates the decision features to properties such as hardness, density, and porosity. Although the output is a material label rather than an image or 3D surface, the experiment demonstrates that NLOS transport can encode interpretable hidden-surface attributes and broadens acoustic NLOS from localization and action recognition toward material-aware environmental perception.

\bookmark[dest=\HyperLocalCurrentHref,level=2]{Robotic Exploration with NLOS Perception}
\subsection{Robotic Exploration with NLOS Perception}
Most classical NLOS experiments focus on reconstructing a pre-defined hidden volume. In robotics, however, hidden perception is valuable because it can actively guide where the robot should move next. Young~\etal~studied single-photon-LiDAR NLOS sensing for autonomous navigation~\cite{youngNavigationNLOS2024}, combining multi-bounce histogram capture, hidden-region occupancy prediction, and a control module that lets a robot choose safer paths in an L-shaped corridor. This work is important because it evaluates NLOS not only by reconstruction fidelity, but by downstream navigation utility. SuperEx~\cite{gargSuperEx2025} later integrates single-photon LiDAR NLOS sensing into indoor mapping and exploration. Instead of treating NLOS reconstruction as a standalone output, it uses timing histograms for NLOS empty-space carving and hidden-structure recovery, then feeds the recovered hidden-space information back into the exploration policy. Together, these works close the loop between NLOS perception and robotic decision making, making NLOS useful even when the reconstruction is partial or uncertain.

\bookmark[dest=\HyperLocalCurrentHref,level=2]{Low-Cost LiDAR NLOS Spatial Reasoning}
\subsection{Low-Cost LiDAR NLOS Spatial Reasoning}
A parallel 2026 trend is the use of consumer or low-cost LiDAR hardware for NLOS reasoning. Somasundaram~\etal~showed that hidden objects can be imaged with consumer LiDAR by using ego-motion to induce a synthetic sampling aperture~\cite{somasundaramConsumerLiDARNLOS2026}. DENALI~\cite{behariDENALI2026} shifts the emphasis from hand-engineered inversion to large-scale data-driven spatial reasoning: it releases 72,000 real-world hidden-object scenes captured as space-time histograms from low-cost LiDARs, enabling models to learn NLOS cues directly from the multi-bounce returns that are usually discarded by consumer depth sensors. Together, these works indicate that mobile NLOS may develop along two complementary paths: physics-based reconstruction with motion-induced sampling, and dataset-driven perception from inexpensive histogram sensors.

\bookmark[dest=\HyperLocalCurrentHref,level=2]{Terahertz NLOS Imaging}
\subsection{Terahertz NLOS Imaging}
Terahertz and sub-terahertz waves provide another route to around-obstacle perception. \href{https://arxiv.org/abs/2205.05066}{Cui and Trichopoulos} showed that building surfaces can act as lossy mirrors for sub-THz signals, enabling hidden-object reconstruction by mirror-folding the propagation geometry around the obstacle. This modality sits between optical and RF NLOS: it has shorter wavelength and potentially higher spatial resolution than many RF systems, while still interacting with common environmental surfaces in a way that can be exploited for around-corner imaging.

\bookmark[dest=\HyperLocalCurrentHref,level=2]{NLOS Human Pose Estimation}
\subsection{NLOS Human Pose Estimation}
Estimating 3D human poses of persons hidden around a corner is an example of semantic NLOS understanding beyond shape reconstruction.

\vspace{0.8mm}
\noindent \textbf{Semantic human-pose recovery from transients.}
Xiao~\etal~move beyond applying a conventional pose estimator to a reconstructed hidden image by directly fusing transient-derived intensity and depth features and regressing semantic body joints end to end~\cite{xiaoNLOSHumanPose2026}. A physics-based pipeline synthesizes training measurements from ordinary smartphone videos, while real experiments demonstrate pose estimation at a 1.75~m relay distance and under an SNR of 0.13. This development advances the earlier physics-based HiddenPose concept toward detailed, low-SNR measured-scene semantics.


\vspace{0.8mm}
\noindent \textbf{Adaptive multi-person pose and mesh sensing.}
Hou~\etal~extend active transient semantics from single-person pose recovery to scenes containing an unknown number of hidden people~\cite{houMultiPersonPose2025}. Their AMPE-NLOS pipeline propagates transient features with LCT to obtain a coarse physical 3D representation, refines it with a 3D U-Net, and predicts both a body-center heatmap and an SMPL parameter map. Body-center-guided sampling then associates mesh parameters with individual people, enabling adaptive multi-person 3D mesh regression. Simulated training data and measurements from a self-built confocal pulsed-laser/SPAD system demonstrate a trajectory in which NLOS reconstruction becomes a front end for structured, multi-instance human understanding rather than only hidden geometry or single-person joints.
 \href{https://arxiv.org/abs/2003.14414}{Isogawa~\etal} introduced an optical NLOS physics-based 3D human-pose-estimation pipeline that converts transient measurements into physically valid body-pose sequences. Although this direction remains less developed than geometric NLOS reconstruction, it illustrates an important path for future systems: hidden-scene perception should ultimately support semantic tasks such as human pose estimation, action recognition, and risk assessment, not only voxelized geometry recovery.